\documentclass[11pt]{article}

\usepackage[letterpaper,margin=1in]{geometry}
\usepackage{microtype}
\usepackage{amsmath}
\usepackage{amssymb}
\usepackage{booktabs}
\usepackage{array}
\usepackage{enumitem}
\usepackage{float}
\usepackage[hidelinks]{hyperref}

\newcommand{\code}[1]{\texttt{#1}}
\newcommand{\rfckw}[1]{\textsc{#1}}

\title{\bfseries A Finer Exploitability Gradation (VLEV)\\ for FedRAMP Rev5 VDR/VER}
\author{Matthew Venne \\ \small Chief Technology Officer, stackArmor}
\date{June 2026}

\begin{document}
\maketitle

\begin{abstract}
\noindent
This is a companion proposal to \emph{A Deterministic, CVSS--Environmental Method for
FedRAMP Rev5 VDR/VER Vulnerability Prioritization}, which derives a Potential Agency
Impact rating (PAIN, N1--N5) and a VDR-TFR-PVR remediation deadline for each finding.
That method treats exploitability as a single yes/no line (Likely-Exploitable or
not). This proposal recommends splitting it into \emph{three} bands so the very
tightest remediation clocks are reserved for the flaws most likely to be exploited in
the wild --- and shows that the change is purely additive to the existing remediation
matrix, slowing nothing down. The key words \rfckw{must}, \rfckw{should}, and
\rfckw{may} are to be read as in RFC~2119.
\end{abstract}

\section{The added inflection point}
The reference method uses one EPSS cut to separate Likely-Exploitable (LEV) from Not
(NLEV). We propose \emph{two} cuts and three bands:
\begin{itemize}[leftmargin=1.5em,nosep,topsep=3pt]
\item \textbf{NLEV} --- $\text{EPSS} \le 0.30$: not likely exploitable.
\item \textbf{LEV} --- $0.30 < \text{EPSS} \le 0.75$: likely exploitable.
\item \textbf{VLEV} --- $\text{EPSS} > 0.75$: \emph{very} likely exploitable.
\end{itemize}
\textbf{Observed active exploitation always maps to VLEV}, regardless of EPSS: if CISA
Vulnrichment reports \code{exploitation = active} (KEV-equivalent), the flaw is being
used \emph{now} and belongs in the most-urgent band by definition.

\section{Alignment with CISA Vulnrichment}
The split is not arbitrary --- it mirrors the taxonomy of the authoritative source.
CISA Vulnrichment's SSVC \emph{Exploitation} metric already distinguishes three
states: \code{none}, \code{poc} (a public proof-of-concept exists), and \code{active}
(observed in the wild). Today's binary LEV/NLEV flattens that distinction; the
proposed bands restore it, with each state lining up to a band and EPSS supplying the
statistical estimate where Vulnrichment is silent:
\begin{itemize}[leftmargin=1.5em,nosep,topsep=3pt]
\item \code{none} $\rightarrow$ NLEV (no known exploit; statistically EPSS $\le 0.30$).
\item \code{poc} $\rightarrow$ LEV (exploit code exists, not yet widespread; EPSS $0.30$--$0.75$).
\item \code{active} $\rightarrow$ VLEV (being used now; EPSS $> 0.75$).
\end{itemize}
As with \code{active}$\rightarrow$VLEV, a known \code{poc} state can floor a finding to
LEV even when its EPSS is lower. The exploitation \emph{state} (when known) and the
EPSS \emph{probability} (always available) then agree on one three-tier scale instead
of being forced into a two-way yes/no.

\section{Calibration: anchored on today's matrix}
We anchor the three bands on FedRAMP's existing deadlines so adoption is trivial and
\emph{nothing remediates slower than today}:
\begin{itemize}[leftmargin=1.5em,nosep,topsep=3pt]
\item \textbf{VLEV reuses FedRAMP's current LEV deadlines, unchanged} --- the
  very-likely / actively-exploited band keeps today's fastest clock.
\item \textbf{LEV} (the new $0.30$--$0.75$ band) gets \emph{slightly more} time than
  VLEV --- likely, but not certain, so a modest grace over the actively-exploited
  band, and a large speed-up versus today, where this band sat in the slow NLEV
  column.
\item \textbf{NLEV} ($\le 0.30$) keeps today's slow deadline as \code{NLEV+NIRV} and
  now \emph{splits} by reachability, so a not-likely-but-internet-reachable flaw
  (\code{NLEV+IRV}) is modestly tighter than a purely internal one.
\end{itemize}
Within every band, internet-reachable (IRV) is tighter than not-reachable (NIRV). The
result is six remediation columns in place of today's three.

\section{The proposed matrix}
Tables~\ref{tab:ab}--\ref{tab:d} give the full six-column grid per Certification Class
(days; $0.5 = 12$ hours). Columns are grouped by reachability; within each group the
three exploitability bands increase left to right. The VLEV columns and
\code{NLEV+NIRV} reproduce the provider's current matrix exactly; \code{LEV} and
\code{NLEV+IRV} are the new, faster columns. PAIN-1 carries no FedRAMP deadline.
Values are illustrative of a representative VDR-TFR-PVR profile; the authoritative day
counts remain FedRAMP's to set.

\begin{table}[H]
\centering
\footnotesize
\renewcommand{\arraystretch}{1.2}
\begin{tabular}{l ccc ccc}
\toprule
& \multicolumn{3}{c}{\textbf{Internet-reachable (IRV)}} & \multicolumn{3}{c}{\textbf{Not reachable (NIRV)}} \\
\cmidrule(lr){2-4}\cmidrule(lr){5-7}
\textbf{PAIN} & VLEV & LEV & NLEV & VLEV & LEV & NLEV \\
\midrule
N5 & 4  & 6   & 16  & 8   & 12  & 32 \\
N4 & 8  & 12  & 48  & 32  & 48  & 64 \\
N3 & 32 & 48  & 96  & 64  & 96  & 192 \\
N2 & 96 & 144 & 176 & 160 & 192 & 192 \\
\bottomrule
\end{tabular}
\caption{Proposed six-column grid --- \textbf{Class A / B} (days; $0.5 = 12$ h).}
\label{tab:ab}
\end{table}

\begin{table}[H]
\centering
\footnotesize
\renewcommand{\arraystretch}{1.2}
\begin{tabular}{l ccc ccc}
\toprule
& \multicolumn{3}{c}{\textbf{Internet-reachable (IRV)}} & \multicolumn{3}{c}{\textbf{Not reachable (NIRV)}} \\
\cmidrule(lr){2-4}\cmidrule(lr){5-7}
\textbf{PAIN} & VLEV & LEV & NLEV & VLEV & LEV & NLEV \\
\midrule
N5 & 2  & 3  & 8   & 4   & 6   & 16 \\
N4 & 4  & 6  & 32  & 8   & 12  & 64 \\
N3 & 16 & 24 & 64  & 32  & 48  & 128 \\
N2 & 48 & 72 & 160 & 128 & 160 & 192 \\
\bottomrule
\end{tabular}
\caption{Proposed six-column grid --- \textbf{Class C} (days; $0.5 = 12$ h).}
\label{tab:c}
\end{table}

\begin{table}[H]
\centering
\footnotesize
\renewcommand{\arraystretch}{1.2}
\begin{tabular}{l ccc ccc}
\toprule
& \multicolumn{3}{c}{\textbf{Internet-reachable (IRV)}} & \multicolumn{3}{c}{\textbf{Not reachable (NIRV)}} \\
\cmidrule(lr){2-4}\cmidrule(lr){5-7}
\textbf{PAIN} & VLEV & LEV & NLEV & VLEV & LEV & NLEV \\
\midrule
N5 & 0.5 & 1  & 4   & 1  & 2   & 8 \\
N4 & 2   & 3  & 16  & 8  & 12  & 32 \\
N3 & 8   & 12 & 48  & 16 & 24  & 64 \\
N2 & 24  & 36 & 160 & 96 & 144 & 192 \\
\bottomrule
\end{tabular}
\caption{Proposed six-column grid --- \textbf{Class D} (days; $0.5 = 12$ h).}
\label{tab:d}
\end{table}

\section{Net effect and adoption}
Nothing slows down. The actively-exploited / $>0.75$ band keeps FedRAMP's current LEV
clock; the broad $0.30$--$0.75$ middle band --- previously parked in the slow NLEV
column --- accelerates onto a near-LEV clock; and reachability now refines the
not-likely band. Because a conforming implementation already drives remediation from a
config-driven \code{[Class][PAIN][column]} lookup, this is additive: it adds columns,
not a redesign. An adopter \rfckw{may} tune the two EPSS cut points ($0.30$, $0.75$)
and the per-Class day counts; the structural recommendation is the three-tier
exploitability axis aligned to CISA Vulnrichment's \code{none}/\code{poc}/\code{active}
states.

\end{document}
